% !TEX root = ../metatime_monograph.tex
\chapter[Sector-Holonomy Release]{Sector-Holonomy Release of the Collatz Quarter-Power Trace}
\label{app:sector_holonomy_release}

\section{Scope and corrected methodological status}

Appendix~\ref{app:collatz_quarter_power} established a comparative signal for
the Collatz-derived exponent $3/4$, but the electron-anchored trace retained a
geometric-mean multiplicative error close to ten.  The sector-holonomy formula
reported here was constructed after inspection of that residual pattern.
Consequently, it is a \emph{retrospective structural candidate}, not a
prospective sealed test.

The corrected status, recorded by the v10.2r1 ledger, is

\[
\boxed{
\text{technical PASS}
\;\land\;
\text{retrospective structural signal}
\;\land\;
\text{physical spectrum FAIL/OPEN}
}
\]

No continuous coefficient was fitted and no particle-dependent residual
correction was introduced.  Nevertheless, the formula-selection history makes
the result ineligible for canonical promotion.  The v10.2r1 classification
supersedes the original v10.2 epistemic label without changing the numerical
operator or its fingerprint.

\section{Base quarter-power coordinate}

Let $A_i$ denote the mass-free Euler--Berry action coordinate and $R_i$ the
mandatory Ramanujan release coordinate.  The frozen v10.1 trace is

\begin{equation}
Q_i=
\frac{(A_e/A_i)^{3/4}-1}{L_3\kappa}
+
\frac{R_i-R_e}{\kappa},
\qquad
\kappa=\frac{\ln2}{24\pi}.
\label{eq:quarter_base_Q}
\end{equation}

The retrospective v10.2 candidate is

\begin{equation}
\ln\frac{m_i}{m_e}=Q_i+G_i+H_i+W_i.
\label{eq:sector_full_operator}
\end{equation}

\section{Collatz generation release}

For the twin-prime pair $(p,p+2)$ associated with an Euler--Berry row, define
its centre $n=p+1$ and ordinary Collatz stopping length $\ell(n)$.  The three
canonical centres satisfy

\[
\ell(4)=2,
\qquad
\ell(6)=8,
\qquad
\ell(12)=9.
\]

Let $c_i$ be the declared colour depth.  The available release depth is
$L_5-c_i$, giving

\begin{equation}
G_i=(L_5-c_i)\frac{\ell_i-\ell_e}{L_3}.
\label{eq:collatz_generation_release}
\end{equation}

For charged leptons, $c_i=0$ and the full intention depth $L_5=5$ is
available.  For quarks, $c_i=3$, leaving the two-dimensional release depth
$L_5-3=L_4=2$.

\section{Signed orientation release}

The pre-mass sector basis supplies a final orientation coordinate $v_{7,i}$
and a chiral path through either $H_+$ or $H_-$.  Define

\[
s_i=
\begin{cases}
+1,& H_+\text{ path},\\
-1,& H_-\text{ path},\\
0,& \text{neutral path}.
\end{cases}
\]

Using the first structural generation $a(i)$ of the same family as a
mass-free reference,

\begin{equation}
H_i=
\frac{s_i}{\max(1,c_i)}
\ln\left(
\frac{1+|v_{7,i}|/\kappa}
     {1+|v_{7,a(i)}|/\kappa}
\right).
\label{eq:signed_orientation_release}
\end{equation}

The colour normalization prevents a colour-depth-three channel from receiving
three independent copies of the same orientation release.

\section{White-Thread down-sector release}

The v3.5 White-Thread artifact defines a mass-free open-holonomy overlap
between the oriented up and down bases.  For each down-type slot $d_i$, the
v10.2 candidate selects the strongest structural channel

\[
w_i=\max_{u_j\in\{u,c,t\}}
\left|\mathcal O_{\mathrm{open}}(u_j,d_i)\right|.
\]

The release term is

\begin{equation}
W_i=L_4\ln w_i
\label{eq:white_thread_release}
\end{equation}

for down quarks and $W_i=0$ otherwise.  The selected channels are
$u\rightarrow d$, $u\rightarrow s$, and $c\rightarrow b$.  The executable
selection uses neither observed CKM elements nor observed masses.  Its choice
as the v10.2 functional form nevertheless belongs to the retrospective
model-selection history and cannot count as independent confirmation.

\section{Numerical retrospective result}

The predictions are

\begin{center}
\begin{tabular}{lrrr}
\toprule
Particle & Quarter-power (GeV) & Sector candidate (GeV) & Validation (GeV)\\
\midrule
$e$   & 0.000510999 & 0.000510999 & 0.000510999\\
$\mu$ & 0.00152075  & 0.108819    & 0.105658\\
$\tau$& 0.0201752   & 1.64836     & 1.77686\\
$d$   & 0.0275743   & 0.00461801  & 0.00467\\
$s$   & 0.121639    & 0.120407    & 0.0934\\
$b$   & 4.17760     & 5.33173     & 4.18\\
$u$   & 0.0522391   & 0.0522391   & 0.00216\\
$c$   & 0.248465    & 1.37826     & 1.27\\
$t$   & 10.3185     & 217.213     & 172.69\\
\bottomrule
\end{tabular}
\end{center}

For the eight non-anchor slots, the mean absolute logarithmic error falls from
$2.2993$ to $0.5137$, and the geometric-mean multiplicative error falls from
$9.967$ to $1.672$.  Excluding only the unresolved $u$ baseline gives

\[
\overline{|\Delta\ln m|}=0.1320,
\qquad
\max|\Delta\ln m|=0.2540,
\qquad
\exp\overline{|\Delta\ln m|}=1.141.
\]

The ordering

\[
e<\mu<\tau,
\qquad
d<s<b,
\qquad u<c<t
\]

is preserved.  These numbers demonstrate that the declared coordinates form a
promising hypothesis generator; they do not constitute an independently
validated derivation.

\section{Remaining obstruction and prospective requirement}

The $u$ prediction remains too high by

\[
\ln\frac{m_u^{\mathrm{pred}}}{m_u^{\mathrm{val}}}=3.1857,
\qquad
\frac{m_u^{\mathrm{pred}}}{m_u^{\mathrm{val}}}=24.18.
\]

This isolates an inter-sector baseline problem.  A $u$-specific subtraction
would be a hidden fit and is forbidden.  Because the charged-fermion masses and
v10.1 residuals have already been inspected, the same table can no longer
serve as a clean prospective holdout for a further operator choice.

The next admissible step is therefore to preregister the allowed structural
inputs, algebraic constraints, and independent validation observable before a
new absolute up-sector baseline formula is evaluated.  Until such an
independent prospective test exists, the sector-holonomy formula remains a
retrospective candidate and Debt~9 remains open.

The numerical executable is
\path{TIR/validation/collatz_sector_holonomy_mass_audit_v10_2.py}.
Its authoritative methodological status is supplied by
\path{TIR/validation/collatz_sector_holonomy_mass_audit_v10_2r1.py}.
